\section{损失、风险、目标与指标不是同一个量}
\label{sec:13-Machine-Learning/01-Learning-Problems-and-Evaluation/02}

\subsection{训练损失掉了八成，线上一动不动}

一家支付公司的风控团队用半年交易记录训练欺诈检测模型。训练日志很好看：交叉熵从 0.693 降到 0.087，验证集准确率停在 94.3\%。模型上线后每天自动标出 200 笔可疑交易交人工复核。第一周结束，运营组反馈：这 200 笔里真欺诈只有 3 笔，同期另有 12 笔欺诈根本没被标出，直接完成了划转。

团队先怀疑优化器。换成 AdamW，加三层残差，训练损失压到 0.042，验证准确率 95.1\%，线上纹丝不动。再怀疑数据不够，训练窗口从半年拉到两年，样本从 80 万涨到 320 万，交叉熵 0.031，准确率 95.8\%，线上仍然没有变化。有人提议换 focal loss，损失掉到 0.018，线上照旧。

到这一步才有人问出那个绕不过去的问题：0.042 这个数究竟测量了什么？它下降说明哪件事变好了？而风控部门在意的拦截率、资金损失、审核工时，又分别是哪些量？

\subsection{四个数字住在四个不同的空间}

答案藏在四个被混用的对象里。在训练脚本里它们相隔几行——\texttt{criterion}、\texttt{loss.backward()}、\texttt{val\_loss}、\texttt{metrics.accuracy}——在数学上它们面对的随机性不同，服务的决定也不同。把它们摞成一架梯子，每一级和上一级之间都隔着一次近似。

\begin{center}
\begin{tikzpicture}[>=stealth, every node/.style={font=\small}]

\node[draw, rounded corners, minimum width=5.6cm, minimum height=0.95cm, align=center] at (0,4.5) {业务后果：拦回的金额 \(-\) 审核成本};
\node[draw, rounded corners, minimum width=5.6cm, minimum height=0.95cm, align=center] at (0,3.0) {部署风险 \(R_{P_{\mathrm{deploy}}}(f)\)};
\node[draw, rounded corners, minimum width=5.6cm, minimum height=0.95cm, align=center] at (0,1.5) {经验风险 \(\widehat R_S(f)\)};
\node[draw, rounded corners, minimum width=5.6cm, minimum height=0.95cm, align=center] at (0,0) {训练目标 \(J_S(\theta)=\widehat R_S+\lambda\Omega\)};

\foreach \y in {4.5,3.0,1.5}
  \draw[<-, thick] (0,\y-0.48) -- (0,\y-1.02);

\node[anchor=east, font=\scriptsize, align=right] at (-2.95,3.75) {代价写得对不对};
\node[anchor=east, font=\scriptsize, align=right] at (-2.95,2.25) {分布偏移};
\node[anchor=east, font=\scriptsize, align=right] at (-2.95,0.75) {代理损失与正则};

\node[draw, dashed, rounded corners, minimum width=2.9cm, minimum height=0.95cm, align=center] at (7.6,4.5) {报告指标};
\draw[->, thick] (2.8,3.2) -- (6.15,4.35);
\node[anchor=west, font=\scriptsize] at (4.0,3.55) {有损压缩};

\end{tikzpicture}
\end{center}

\emph{优化器只在最底下那一级工作，业务后果在最上面，中间隔着三次可能出错的翻译。}

风控团队做的所有努力——换优化器、加数据、换损失——都发生在梯子最底下那一级。那一级确实一直在改善。

\subsection{损失是一次预测的一次记账}

对单个样本，预测 \(f(x)\) 与真实结果 \(y\) 之间的代价写作 \(\ell(f(x),y)\)。平方损失 \((f(x)-y)^2\) 对大误差二次增长，残差从 2 变到 4，代价翻的是四倍；绝对损失 \(\abs{f(x)-y}\) 与距离成正比，不给离群点加权；对数损失 \(-\log p_{\mathrm{model}}(y\given x)\) 按模型分给真实结果的概率计费，分了 0.01 就记 4.6，分了 0.9 只记 0.105。

损失的形状决定训练时哪类错误被优先修。而同一组 \((x,y)\) 上可以写出好几个合理的损失，它们回答的问题并不相同。若 \(Y\given X=x\) 的分布偏斜，平方损失的总体最优预测是条件均值 \(\E[Y\given X=x]\)，绝对损失的总体最优预测是条件中位数。两者都没算错，只是在估计不同的统计量。

一个离数学不远的场景可以把这件事钉死。便利店每天下午五点前要定次日的面包订量 \(q\)，需求 \(D\) 尚未发生。每剩一条报废成本 \(c_{\mathrm o}=4\) 元，每缺一条折合流失成本 \(c_{\mathrm u}=16\) 元。给定天气、星期和周边活动信息 \(X=x\)，条件需求右偏：多数日子卖 80 到 120 条，碰上球赛能冲到 200 条以上。订量 \(q\) 的条件风险是

\[
C(q\given x)=c_{\mathrm o}\E[(q-D)_+\given X=x]+c_{\mathrm u}\E[(D-q)_+\given X=x].
\]

若条件分布函数 \(F(\cdot\given x)\) 连续，求导并令其为零，

\[
\begin{aligned}
\frac{\partial C(q\given x)}{\partial q}
&=c_{\mathrm o}F(q\given x)-c_{\mathrm u}\bigl(1-F(q\given x)\bigr)=0,\\[2pt]
F(q^\star\given x)&=\frac{c_{\mathrm u}}{c_{\mathrm u}+c_{\mathrm o}}=\frac{16}{16+4}=0.8 .
\end{aligned}
\]

最优订量是条件分布的第 80 百分位数，不是条件均值。缺货越贵，所需分位数越高。用均方误差训出来的模型学的是均值，拿它去订货，总成本会比最优水平高出一截，具体多少取决于分布的偏斜程度。损失函数不是在调参数，它在指定模型该学哪个统计量。

\subsection{风险是把代价摊到未来分布上}

真正想控制的是预测器在所有未来场景上的期望表现。若部署数据来自 \(P\)，

\[
R_P(f)=\E_{(X,Y)\sim P}\bigl[\ell(f(X),Y)\bigr].
\]

这个量算不出来，因为 \(P\) 未知。手头只有 \(S=\set{(x_i,y_i)}_{i=1}^n\)，于是拿经验风险顶上：

\[
\widehat R_S(f)=\frac1n\sum_{i=1}^n\ell(f(x_i),y_i).
\]

对一个在见到数据之前就固定的 \(f\)，样本独立同分布且损失可积时，经验风险是总体风险的无偏估计，大数定律保证它收敛到 \(R_P(f)\)。

但训练不是评估一个预先固定的 \(f\)，训练是从函数类里挑出让经验风险最小的 \(\widehat f\)。挑选本身要消耗样本里的随机性。设想 100 个候选函数的真实风险完全相同，只因训练数据的抽样波动，总有一个碰巧在这份数据上最好；你把它挑出来，它并不泛化更好，只是在这份噪声上更走运。于是 \(\widehat R_S(\widehat f)\) 系统性地低于 \(R_P(\widehat f)\)。同一份噪声被用了两次：一次用来选函数，一次用来给这个函数打分。

风险还必须标明下标。若 \(P_{\mathrm{train}}\ne P_{\mathrm{deploy}}\)，样本量趋于无穷也只让经验风险逼近 \(R_{P_{\mathrm{train}}}(f)\)，不会自动逼近 \(R_{P_{\mathrm{deploy}}}(f)\)。图上那一级台阶不会因为数据变多而消失。

便利店的数据库里还埋着更隐蔽的一层。记录下来的是销量 \(S\) 而非需求 \(D\)，库存为 \(I+q\) 时 \(S=\min\set{D,\,I+q}\)。售罄日只告诉你需求至少到了库存量，没告诉你超出多少。把销量当需求训练，模型会系统性低估旺盛需求，于是少订货、制造更多售罄记录、进一步压低未来的``需求''数据。这不是随机波动，是历史行动决定了标签能被观察到多少。

\subsection{训练目标在风险之上又加了两层近似}

实际跑优化时，目标函数通常写成

\[
J_S(\theta)=\widehat R_S(f_\theta)+\lambda\Omega(\theta).
\]

第一项拟合数据，第二项通过范数惩罚、稀疏约束或别的结构限制缩小候选范围。\(\Omega(\theta)\) 可以表达对平滑的偏好、对稀疏的偏好、对群组结构的偏好，或同时承担几件事。它改善的是从有限样本里选出稳定函数的概率，而不是让函数在总体上更正确——它用偏差换方差，前提是这笔交换在未见数据上确实换来更小的期望误差。

第二层近似是代理损失。分类真正关心的常是 \(0\)-\(1\) 错误 \(\ell_{01}(f(x),y)=\ind\set{f(x)\ne y}\)，而它分段常数、梯度几乎处处为零，优化器不知道往哪儿动。于是换成 logistic loss、交叉熵或 hinge loss：光滑、凸、可导，局部梯度给出明确方向。副作用各有各的：logistic loss 对已分对的样本仍施加非零惩罚，鼓励模型不但分对还要分得有把握；hinge loss 对离边界足够远的正确样本零梯度，分对了就不再推。

这些替代损失在一定条件下与 \(0\)-\(1\) 损失有相同的总体最优解，但一致性需要条件：函数类足够灵活、数据量足够大、噪声结构与替代损失的内在假设相容。类别严重不平衡、代价不对称或决策阈值不取 \(1/2\) 时，替代损失的最优解可能偏离你真正关心的边界。选它不能只因为``方便优化''。

至于优化算法，SGD、Newton、坐标下降、Adam 解决的是同一个 \(J_S\) 上的不同搜索策略。把 \(J_S\) 优化到机器精度，修不了损失定义写错、数据泄漏或分布选错，它只是忠实完成了交给它的数值任务。

\subsection{指标是一次有损压缩}

评价指标把一堆预测压成几个便于比较的数。准确率、精确率、召回率、\(F_1\)、ROC AUC、均方误差、校准误差各自保留和丢弃不同的信息。一个指标改善，不意味着别的指标改善，更不意味着部署效果改善。

类别极不平衡时这件事最刺眼。若 98\% 的交易正常，恒预测``正常''就有 98\% 准确率，比多数精心训练的模型都高，而它从未标记过一笔欺诈，召回率为零，风控价值为零。准确率把两类错误等权相加，比例悬殊时它完全被多数类主导。把准确率从 98.2\% 提到 98.7\%，看着是 0.5 个百分点的进步，这 0.5 个百分点可能全部来自更多正常交易被正确放行，欺诈那一侧毫无变化甚至退化。

AUC 衡量的是随机抽一正一负、模型给正例更高分的概率。它不依赖阈值，也不受类别先验影响，代价是它不告诉你别的。AUC 为 0.95 的模型可以把所有正例排在 0.51 到 0.55 之间、所有负例排在 0.50 以下——排序对，间距细如发丝，一点扰动就翻。它也不告诉你排序顶部的精确率是多少，而每天只复核 200 笔时，你要的恰恰是最可疑那 200 笔的精确率。

\(F_1\) 是精确率与召回率的调和平均。用调和平均意味着两者都高才高：精确率 0.9、召回率 0.1 时 \(F_1\) 只有 0.18，而不是算术平均的 0.5。\(F_1\) 完全忽略真阴性，正常交易被正确放行不改善它。在欺诈检测里这或许合理，在疾病筛查里就未必——正确排除一位健康人是有价值的决策。选 \(F_1\) 本身就是一次立场声明。

\subsection{总体指标可能与每个分组的指标方向相反}

汇总还会掩盖分布。总体准确率从 94\% 升到 95\%，同时新注册用户那一组从 88\% 掉到 82\%，涨幅全部来自老用户，而新用户的体验在恶化——这在汇总数字里看不见。微平均按样本等权加总，结果被大群体主导；宏平均先算每组指标再平均，赋予各组等权。选哪个，取决于你认为该等权的是样本还是群体，而群体规模悬殊时这两条原则给出相反的方向。

更极端的一种情形是方向整体翻转：模型 A 在每一个分组上都不如模型 B，合并之后 A 的总体指标反而更高。这不需要任何统计错误，只要两组的样本量与基础难度配比合适就会发生，它与\hyperref[sec:13-Machine-Learning/12-Causal-Inference/03]{``混杂、碰撞点与后门调整''}中讨论的 Simpson 反转是同一个结构。汇总是对分组做加权平均，而权重由分组人数决定；当两个模型面对的分组人数不同——例如 A 的流量里困难分组占比更小——加权平均就可能与逐组比较相悖。所以模型对比表里除了总体一行，还必须有分组几行；只报总体的对比，读者无法排除这种反转。

\subsection{阈值由代价决定，不由概率模型自动给出}

设 \(p(x)=P(Y=1\given X=x)\)。判为正类而真实为负，承担假阳性代价 \(C_{\mathrm{FP}}\)；判为负类而真实为正，承担假阴性代价 \(C_{\mathrm{FN}}\)。给定 \(x\)，两种行动的条件风险是

\[
\begin{aligned}
\rho(1\given x)&=C_{\mathrm{FP}}\bigl(1-p(x)\bigr),\\
\rho(0\given x)&=C_{\mathrm{FN}}\,p(x).
\end{aligned}
\]

选正类当且仅当 \(\rho(1\given x)\le\rho(0\given x)\)，即

\[
p(x)\ge\frac{C_{\mathrm{FP}}}{C_{\mathrm{FP}}+C_{\mathrm{FN}}} .
\]

只有两类代价相等时阈值才是 \(1/2\)。漏掉一笔欺诈的代价若是错误冻结一笔正常交易的 50 倍，阈值就落到 \(1/51\approx 0.02\) 附近，而不是 0.5。精确率与召回率的取舍就在这个式子里：调阈值即沿着这条取舍曲线滑动，落点由代价比决定，不由模型决定。概率模型的任务是输出校准良好的 \(p(x)\)，决策层再结合代价定阈值；若模型给的只是排序分数而非校准概率，这个公式根本不能直接用。这条接口的完整展开见\hyperref[ch:092]{``统计决策与可信预测：模型给出概率之后怎样行动''}。

代价还可能随 \(x\) 变，也可能撞上容量约束。医院每天只有 40 个复查名额时，决策变成带总预算的分配问题，不能再对每位患者独立阈值化，而要按风险排序取前 40 名。

\subsection{没有基线的比较不构成证据}

回到梯子。所有这些数字——0.042、95.8\%、AUC 0.95——单独看都不成立，因为它们没有参照物。一个指标值本身不说明任何事，说明问题的是它与一条不含模型的规则之间的差距。

该有的基线不止一条。最朴素的是常数基线：分类任务恒输出多数类，回归任务恒输出训练集均值，它给出准确率和均方误差的地板；98\% 准确率若地板就在 98\%，这个模型没有产出任何信息。第二条是随机基线：按类别先验随机猜，它给出 AUC 的地板 0.5，也给出精确率的地板等于正例率。第三条是持续性基线，在时间序列和用户行为预测里最狠——用``昨天的值''或``用户上次的选择''直接当预测，很多号称精妙的模型跑不赢它。第四条是单特征基线：只用最强的那个特征做逻辑回归或单变量阈值；若三亿参数的模型只比它高一个百分点，多出来的复杂度换来了什么需要另行解释。第五条是当前线上规则，也就是模型要替换掉的那套人工策略，这是唯一与决策直接可比的基线。

基线还必须与模型走完全相同的评估流程：同样的划分、同样的预处理、同样的指标口径。基线用全量数据的均值、模型用训练集的均值，比较就失去意义。选定基线后，差距还要放在抽样波动里看：测试集上 0.3 个百分点的领先，若置信区间宽达 2 个百分点，那不是领先，是噪声。区间怎么算、划分怎么做，是下一节和本单元第四篇的事。

\subsection{回到那四行日志}

现在再看风控团队的处境。交叉熵从 0.693 降到 0.042，说明优化器在代理损失上干得很好；验证准确率 94.3\%，说明模型在等权计错的意义上不差。而真实问题在梯子的每一级上都错了一点。

交叉熵鼓励输出校准概率，却不区分两类错误的业务代价——误标一笔正常交易（浪费人工）和漏掉一笔欺诈（资金损失）在它眼里是对称的。94.3\% 的准确率里有约 96 个百分点来自正确放行的正常交易，模型完全可能只抓到两成欺诈而准确率照样好看。训练目标里的 L2 正则在把概率往 0.5 拉，对本应接近 0 或 1 的极端样本，这种收缩直接削掉了区分能力。而运营每天只能复核 200 笔这条容量约束，从头到尾没有进入任何一个损失函数或指标。

所以看到目标下降而部署效果不动时，要分四层追问：样本损失惩罚的是哪类错误，它对应的最优预测是均值、中位数、分位数还是概率，那是你要模型学的统计量吗；风险是对哪个分布取期望，训练样本来自那个分布吗，标签有没有被历史行动截断；训练目标额外加了什么正则或代理，它改的是约束还是问题本身；报告指标丢掉了什么，在代价不对称、群体有差异、存在容量约束时，它还能不能为行动提供正确的排序。

这四层在一个 Python 文件里可能只隔几行，在数学上是四个对象，修好一层不会自动修好另外三层。下一节把梯子最底下那一级拆开：训练目标里的那个 \(\ell\) 并不是随手选的，它背后站着一个关于噪声的假设。

